% GENERATED FILE. Edit canonical lessons, then run npm run generate:books.
% canonical-source-hash: fnv1a64:a27df012d8ddebe3
% canonical-lessons: ML-C203-ruciyulla, ML-C203-cinna, ML-C203-ayanna, ML-C203-irukiya, ML-C203-kaliyaya

\chapter{Tasty, Rotten, Loose, Tight, Empty}
\label{ch:ml-tasty-rotten-loose-tight-empty}
\hlchaptermodality{203}

\begin{chapteropening}
\textbf{By the end of this chapter:} \emph{I can say tasty, rotten, loose, tight and empty.}

Complete the last lesson of chapter 203: I can say tasty, rotten, loose, tight and empty.
\end{chapteropening}

\section[\texorpdfstring{\ml{രുചിയുള്ള} (ruciyuḷḷa)}{രുചിയുള്ള (ruciyuḷḷa)}]{\ml{രുചിയുള്ള} (ruciyuḷḷa) — tasty}

\label{lesson:ML-C203-ruciyulla}



\begin{quote}
Before the new one: say the Malayalam for a drawback, then the Malayalam for fairness.
\end{quote}

\subsection*{You'll want to know: \ml{രുചിയുള്ള}}
\textbf{\ml{രുചിയുള്ള}} — \emph{ruciyuḷḷa} — \textquotedblleft{}tasty\textquotedblright{}.

\begin{grammarlens}[title={What you've built}]
One more word for this chapter.
\end{grammarlens}

\subsection*{Your turn}
\begin{itemize}
\raggedright
  \item \emph{Say it:} \emph{ruciyuḷḷa}
  \item \emph{Say it:} \emph{ruciyuḷḷa}, once more
  \item \emph{Say it:} say \emph{ruciyuḷḷa}
  \item \emph{Recall:} say the Malayalam for a drawback, then the Malayalam for fairness, then say \emph{ruciyuḷḷa} again
\end{itemize}

\begin{tcolorbox}[breakable,colback=teal!4,colframe=teal!35!black,title={Before you move on}]
What does \ml{രുചിയുള്ള} mean? (\textquotedblleft{}Tasty\textquotedblright{}.) Say it once more.
\end{tcolorbox}
\section[\texorpdfstring{\ml{ചീഞ്ഞ} (cīñña)}{ചീഞ്ഞ (cīñña)}]{\ml{ചീഞ്ഞ} (cīñña) — rotten}

\label{lesson:ML-C203-cinna}



\begin{quote}
Before the new one: say the Malayalam for fairness, then the Malayalam for tasty.
\end{quote}

\subsection*{You'll want to know: \ml{ചീഞ്ഞ}}
\textbf{\ml{ചീഞ്ഞ}} — \emph{cīñña} — \textquotedblleft{}rotten\textquotedblright{}.

\begin{grammarlens}[title={What you've built}]
One more word for this chapter.
\end{grammarlens}

\subsection*{Your turn}
\begin{itemize}
\raggedright
  \item \emph{Say it:} \emph{cīñña}
  \item \emph{Say it:} \emph{cīñña}, once more
  \item \emph{Say it:} say \emph{cīñña}
  \item \emph{Recall:} say the Malayalam for fairness, then the Malayalam for tasty, then say \emph{cīñña} again
\end{itemize}

\begin{tcolorbox}[breakable,colback=teal!4,colframe=teal!35!black,title={Before you move on}]
What does \ml{ചീഞ്ഞ} mean? (\textquotedblleft{}Rotten\textquotedblright{}.) Say it once more.
\end{tcolorbox}
\section[\texorpdfstring{\ml{അയഞ്ഞ} (ayañña)}{അയഞ്ഞ (ayañña)}]{\ml{അയഞ്ഞ} (ayañña) — loose}

\label{lesson:ML-C203-ayanna}



\begin{quote}
Before the new one: say the Malayalam for tasty, then the Malayalam for rotten.
\end{quote}

\subsection*{You'll want to know: \ml{അയഞ്ഞ}}
\textbf{\ml{അയഞ്ഞ}} — \emph{ayañña} — \textquotedblleft{}loose\textquotedblright{}.

\begin{grammarlens}[title={What you've built}]
One more word for this chapter.
\end{grammarlens}

\subsection*{Your turn}
\begin{itemize}
\raggedright
  \item \emph{Say it:} \emph{ayañña}
  \item \emph{Say it:} \emph{ayañña}, once more
  \item \emph{Say it:} say \emph{ayañña}
  \item \emph{Recall:} say the Malayalam for tasty, then the Malayalam for rotten, then say \emph{ayañña} again
\end{itemize}

\begin{tcolorbox}[breakable,colback=teal!4,colframe=teal!35!black,title={Before you move on}]
What does \ml{അയഞ്ഞ} mean? (\textquotedblleft{}Loose\textquotedblright{}.) Say it once more.
\end{tcolorbox}
\section[\texorpdfstring{\ml{ഇറുകിയ} (iṟukiya)}{ഇറുകിയ (iṟukiya)}]{\ml{ഇറുകിയ} (iṟukiya) — tight}

\label{lesson:ML-C203-irukiya}



\begin{quote}
Before the new one: say the Malayalam for rotten, then the Malayalam for loose.
\end{quote}

\subsection*{You'll want to know: \ml{ഇറുകിയ}}
\textbf{\ml{ഇറുകിയ}} — \emph{iṟukiya} — \textquotedblleft{}tight\textquotedblright{}.

\begin{grammarlens}[title={What you've built}]
One more word for this chapter.
\end{grammarlens}

\subsection*{Your turn}
\begin{itemize}
\raggedright
  \item \emph{Say it:} \emph{iṟukiya}
  \item \emph{Say it:} \emph{iṟukiya}, once more
  \item \emph{Say it:} say \emph{iṟukiya}
  \item \emph{Recall:} say the Malayalam for rotten, then the Malayalam for loose, then say \emph{iṟukiya} again
\end{itemize}

\begin{tcolorbox}[breakable,colback=teal!4,colframe=teal!35!black,title={Before you move on}]
What does \ml{ഇറുകിയ} mean? (\textquotedblleft{}Tight\textquotedblright{}.) Say it once more.
\end{tcolorbox}
\section[\texorpdfstring{\ml{കാലിയായ} (kāliyāya)}{കാലിയായ (kāliyāya)}]{\ml{കാലിയായ} (kāliyāya) — empty}

\label{lesson:ML-C203-kaliyaya}



\begin{quote}
Before the new one: say the Malayalam for loose, then the Malayalam for tight.
\end{quote}

\subsection*{You'll want to know: \ml{കാലിയായ}}
\textbf{\ml{കാലിയായ}} — \emph{kāliyāya} — \textquotedblleft{}empty\textquotedblright{}.

\begin{grammarlens}[title={What you've built}]
One more word for this chapter.
\end{grammarlens}

\subsection*{Your turn}
\begin{itemize}
\raggedright
  \item \emph{Say it:} \emph{kāliyāya}
  \item \emph{Say it:} \emph{kāliyāya}, once more
  \item \emph{Say it:} say \emph{kāliyāya}
  \item \emph{Recall:} say the Malayalam for loose, then the Malayalam for tight, then say \emph{kāliyāya} again
\end{itemize}

\begin{tcolorbox}[breakable,colback=teal!4,colframe=teal!35!black,title={Before you move on}]
What does \ml{കാലിയായ} mean? (\textquotedblleft{}Empty\textquotedblright{}.) Say it once more.
\end{tcolorbox}
